\section{Research Objectives}
\label{sec:objectives}

The primary goal of this research is to develop a mathematically rigorous, multi--scale computational framework that derives macroscopic environmental transport parameters directly from low--Reynolds--number microstructural fluid--particle mechanics. To achieve this overarching goal, the project will fulfill the following specific objectives:

\begin{enumerate}
    \item \textbf{Implement a High--Fidelity Stokesian Dynamics Module:} Configure and optimize an open--source multi--body hydrodynamic tracking engine (e.g., LAMMPS or specialized SD libraries) to simulate the trajectories of rigid particulate suspensions ($Re \ll 1$) within complex benthic geometries. This includes explicitly incorporating far--field many--body mobility matrices, short--range lubrication forces, and the ambient fluid's rate--of--strain tensor ($\bm{{E}}^{\infty}$) under localized shear conditions.
    \item \textbf{Characterize Microstructural Particulate Trapping Mechanisms:} Quantify the discrete capture efficiency and accumulation dynamics of colloidal particles interacting with complex boundary interfaces, specifically mapping the physical constraints of coral canopy architectures (e.g., boundary layer interception, wake--vortex stagnation) and high--permeability sediment pore throats characteristic of the Tayrona National Natural Park bays \cite{goldman1967shear, sen2010lightdriven}.
    \item \textbf{Develop a Deterministic Micro--to--Macro Upscaling Protocol:} Establish a mathematical mapping interface to upscale discrete particle configuration data into continuous fields. This objective entails calculating the Effective Hydrodynamic Dispersion Tensor ($\bm{D}_{\text{eff}}$) via long--time asymptotic particle displacement variance, and calculating the continuous kinetic trapping rate ($R$) via net cross--boundary mass fluxes.
    \item \textbf{Validate Against Empirical Experimental Literature:} Execute high--order numerical solvers for the continuum Advection--Diffusion--Reaction (ADR) equation using the deterministically derived parameters. Critically validate the predictive accuracy of the upscaled computational framework by directly contrasting the numerical concentration fields and trapping efficiencies against independent, empirical datasets available in the literature, including physical coral canopy flume tracking data and local reef sand permeability measurements \cite{mendrik2024microplastic, reidenbach2006boundary}.
\end{enumerate}
